%; whizzy chapter -dvi
% -initex iniptex -latex platex -format platex -bibtex jbibtex -fmt fmt
% 以上 whizzytex を使用する場合の設定。

%     Tokyo Debian Meeting resources
%     Copyright (C) 2012 Junichi Uekawa
%     Copyright (C) 2011, 2015 Nobuhiro Iwamatsu

%     This program is free software; you can redistribute it and/or modify
%     it under the terms of the GNU General Public License as published by
%     the Free Software Foundation; either version 2 of the License, or
%     (at your option) any later version.

%     This program is distributed in the hope that it will be useful,
%     but WITHOUT ANY WARRANTY; without even the implied warranty of
%     MERCHANTABILITY or FITNESS FOR A PARTICULAR PURPOSE.  See the
%     GNU General Public License for more details.

%     You should have received a copy of the GNU General Public License
%     along with this program; if not, write to the Free Software
%     Foundation, Inc., 51 Franklin St, Fifth Floor, Boston, MA  02110-1301 USA

%  preview (shell-command (concat "evince " (replace-regexp-in-string "tex$" "pdf"(buffer-file-name)) "&"))

%%ここからヘッダ開始。

\documentclass[mingoth,a4paper]{jsarticle}
\usepackage{monthlyreport}
% 日付を定義する、毎月変わります。
\newcommand{\debmtgyear}{2019}
\newcommand{\debmtgmonth}{12}
\newcommand{\debmtgdate}{22}
% started from zero:
% (let ((year 2013) (month 7)) (+ (* (- year 2005) 12) month -1))
\newcommand{\debmtgnumber}{181}

% Needed to import pandoc-generated LaTeX documents.
% See https://stackoverflow.com/questions/40438037/tightlist-error-using-pandoc-with-markdown
\providecommand{\tightlist}{%
  \setlength{\itemsep}{0pt}\setlength{\parskip}{0pt}}

% tikz picture の為のマクロ設定
\usepackage[dvipdfmx]{graphicx}
\usepackage{tikz}

\begin{document}

\begin{titlepage}
\thispagestyle{empty}
% タイトルページ:編集必要な部分は最初のマクロに飛ばすこと

\vspace*{-2cm}
第\debmtgnumber{}回 東京エリア Debian 勉強会資料\\
\hspace*{-2cm}
\includegraphics{image2012-natsu/dotdeb.pdf}\\
\hfill{}\debmtgyear{}年\debmtgmonth{}月\debmtgdate{}日

% ここはアップデートすること
% 全角文字にしないとフォントのサイズが合わないので注意
\rotatebox{10}{\fontsize{30}{30} {\gt　ｎｆｔａｂｌｅｓ入門}}\\

\vspace*{-2cm}
\hfill{}\includegraphics[height=6cm]{image200502/openlogo-nd.eps}
\end{titlepage}

\newpage

\begin{minipage}[b]{0.2\hsize}
 \definecolor{titleback}{gray}{0.9}
 \colorbox{titleback}{\rotatebox{90}{\fontsize{80}{80} {\gt デビアン勉強会} }}
\end{minipage}
\begin{minipage}[b]{0.8\hsize}
\hrule
\vspace{2mm}
\hrule
\begin{multicols}{2}
\tableofcontents
\end{multicols}
\vspace{2mm}
\hrule
\end{minipage}

\dancersection{最近のDebian関連のミーティング報告}{杉本典充}

\subsection{2019年10月度 東京エリアDebian勉強会}

2019年10月19日(土)に東京エリアDebian勉強会を開催しました。会場は荒川区立町屋文化センターの会議室を借りて行いました。参加者は5名でした。

セミナーは杉本さんの「Debian GNU/kFreeBSD セットアップガイド 2019年版」を行いました。Debian GNU/kFreeBSDがDebian Portsに移行となりインストール方法や使えるパッケージが少ないながらどこまで動くのか試してみた内容でした。

「月間 Debian Policy」の時間は、最近更新があったDebian Policyはどのような課題があって変更に至ったのか、勉強会の参加している人で意見を出し合い共有しました。

% 勉強会の後は参加者で懇親を深めました。

\subsection{OSC 2019 Tokyo/Fall}

2019年11月23日(土)にDebian JP Project／東京エリアDebian勉強会は、OSC 2019 Tokyo/Fall\footnote{\url{https://www.ospn.jp/osc2019-fall/}}に参加しました。OSCの会場は明星大学様でした。

OSCのイベント全体では、11/23(土)に約500名、11/24(日)に約330名の方が来場しました。

セミナーでは「Debian Updates」という表題でDebian 10 busterのリリース情報の説明を杉本が発表し、参加者は19名でした。またブース展示を行い、約43名のイベント参加者と交流しました。

% 勉強会の後は参加者で懇親を深めました。


\dancersection{事前課題}{杉本典充}

今回の事前課題は以下です。

\begin{enumerate}
\item nftablesを使ったことはありますか。OSは問いません。
\item Debianの話で聞きたいことがあれば教えてください。
\end{enumerate}

%この課題に対して提出いただいた内容は以下です。

\begin{multicols}{2}
{\small
\begin{prework}{ koedoyoshida }
  \begin{enumerate}
  \item $B;H$C$?$3$H$O$"$j$^$;$s(B
  \item ($B2sEz$J$7(B)
  \end{enumerate}
\end{prework}

\begin{prework}{ NOKUBI Takatsugu (knok) }
  \begin{enumerate}
  \item $B;H$C$?$3$H$O$"$j$^$;$s(B
  \item Cosmo Communicator$B$G(BDebian$B$,F0$/$=$&$G$9$,C/$+Gc$$$^$9$+(B?
  \end{enumerate}
\end{prework}

\begin{prework}{ Kouhei Maeda (mkouhei) }
  \begin{enumerate}
  \item $B;H$C$?$3$H$O$"$j$^$;$s(B
  \item ($B2sEz$J$7(B)
  \end{enumerate}
\end{prework}

\begin{prework}{ yy\_y\_ja\_jp }
  \begin{enumerate}
  \item $B;H$C$?$3$H$O$"$j$^$;$s(B
  \item ($B2sEz$J$7(B)
  \end{enumerate}
\end{prework}

\begin{prework}{ dictoss }
  \begin{enumerate}
  \item $B;H$C$?$3$H$,$"$j$^$9(B
  \item $B%\!<%I(BPC$B7O$N>pJs!J(BRaspberry Pi$B!"(BPINE64$B!K!"%3%s%F%J$N%2%9%H(BOS$B$H$7$F(Bdebian$B$O$I$NDxEY8~$$$F$$$k$N$+$ND4::!&8+2r(B
  \end{enumerate}
\end{prework}

}
\end{multicols}

%\dancersection{Debian Trivia Quiz}{username}
%
%Debianの昨今の話題についてのQuizです。
%
%今回の出題範囲は\url{debian-devel-announce@lists.debian.org} や \url{debian-news@lists.debian.org}などに投稿された内容からです。
%
%\begin{multicols}{2}
%%; whizzy-master ../debianmeetingresume201211.tex
% $B0J>e$N@_Dj$r$7$F$$$k$?$a!"$3$N%U%!%$%k$G(B M-x whizzytex $B$9$k$H!"(Bwhizzytex$B$,MxMQ$G$-$^$9!#(B
%

\santaku
{DebConf13 $B$N3+:ECO$H3+:EF|$O!)(B}
{$BF|K\!"El5~ET(B 6$B7n(B20$BF|(B}
{$B%K%+%i%0%"(B $B%^%J%0%"(B 7$B7n(B8-14$BF|(B}
{$B%9%$%9!"%t%)!<%^%k%-%e(B 8$B7n(B11-18$BF|(B}
{3}
{$B%K%+%i%0%"$O(BDebConf12$B$N3+:ECO$G$9!#(B
DebConf13$B$O%9%$%9$N%-%c%s%WCO$G3+:E$G$9!#(B
6/20$B$O3'$5$sM=Dj$r6u$1$F$*$-$^$7$g$&!#(B}

\santaku
{$B@$3&$N(BWeb$B%5!<%P$G:G$b?M5$$N$"$k(BLinux $B%G%#%9%H%j%S%e!<%7%g%s(B(W3Techs$BD4$Y(B)$B$O!)(B}
{CentOS}
{Debian}
{Ubuntu}
{B}
{\url{http://w3techs.com/technologies/history_details/os-linux}$B$K7k2L$N%0%i%U$,$"$j$^$9!#(B
$B8=:_(B Linux $B$r;HMQ$7$F$$$k(B web $B%5!<%P$N(B 32.9\% $B$,(B Debian $B$rMxMQ$7$F$*$j!"$=$N3d9g$O8=:_$bA}2C$rB3$1$F$$$k$=$&$G$9!#(B}

\santaku
{Debian $B%+!<%M%k%A!<%`$N%a%s%P!<$G$"$j!"(Bkernel.org $B$N(B 3.2.y $B0BDjHG7ONs$N%a%s%F%J$G$b$"$k(B Ben Hutchings $B$5$s$,<!4|(B Debian $B0BDjHG$H0l=o$K=P2Y$5$l$k(B Linux $B%+!<%M%k$K(B (3.2 $B7ONs$N(B mainline $B$K$OL5$$(B) $BDI2C5!G=$,Ek:\$5$l$kM=Dj$G$"$k$H=R$Y$F$$$^$9!#(B
$BB?$/$NDI2CE@$NCf$K4^$^$l$J$$$b$N$O2?!)(B}
{PREEMPT\_RT}
{Hyper-V guest drivers$B$N6/2=(B}
{ARM64/AArch64$B%"!<%-%F%/%A%c%5%]!<%H(B}
{C}
{Hyper-V guest drivers$B$O(Bmainline kernel$B$G(B3.2$B$K$b4^$^$l$F$$$^$9$,!"$h$j2~A1$5$l$?(B3.4$B$+$i$N=$@5$,F3F~$5$l$^$9!#(B
PREEMPT\_RT$B$O%O!<%I%j%"%k%?%$%`$r<B8=$9$k$?$a$N(BPatch$B!"(B
linux-image-rt-amd64 , linux-image-rt-686-pae $B$N(Bmetapackage$B$G;HMQ$G$-$^$9!#(B
$B?7$7$$(BARM 64$B%S%C%H%"!<%-%F%/%A%c%5%]!<%H$O(Bmainline kernel 3.7$B$+$i(B}

\santaku
{Wookey$B$5$s$,%"%J%&%s%9$7$?(Balpha$BHG$N(BDebian port arm64 image$B$O!)(B}
{Debian/Ubuntu port image}
{Debian/KFreeBSD port image}
{Debian/GnuHurd port image}
{A}
{self-bootstrapp(non x86)$BBP1~$H$N$3$H$G$9!#(B\url{http://wiki.debian.org/Arm64Port}$B$G%9%F!<%?%9$,3NG'$G$-$^$9!#(B}

\santaku
{700,000$BHVL\$N%P%0$,Js9p$5$l$?F|$rEv$F$k(B700000thBugContest$B$N7k2L$,=P$^$7$?!#$=$NM=A[F|$HJs9pF|$O!)(B}
{2012/12/12$B$rM=A[$7$?(BDavidPrevot}
{$BM=A[F|(B:2013/02/04$B!"Js9pF|(B:2013/02/14}
{$BM=A[F|(B:2013/02/07$B!"Js9pF|(B:2013/02/14}
{$BM=A[F|(B:2013/02/14$B!"Js9pF|(B:2013/02/07}
{C}
{$B:G$b6a$$(B2013/02/14$B$rM=A[$7$?(BChristian Perrier$B$5$s$,Ev$F$^$7$?!#7k2L$O(B\url{http://wiki.debian.org/700000thBugContest}$B$G8x3+$5$l$F$$$^$9!#(B
$B$^$?!"(B800,000/1,000,000$BHVL\$N%P%0$,Js9p$5$l$kF|$rEv$F$k%3%s%F%9%H(B\url{http://wiki.debian.org/800000thBugContest}$B$b3+:E$5$l$F$$$^$9!#(B}

\santaku
{master.debian.org$B$,?7$7$$5!3#$K0\9T$5$l$^$7$?!#$3$l$O2?$N%5!<%P$G$7$g$&$+(B $B!)(B}
{@debian.org$B$N%a!<%k%5!<%P(B}
{$B%Q%C%1!<%8$N%^%9%?!<%5!<%P(B}
{$B%Q%C%1!<%8$N%9%]%s%5!<(B(mentor)$B$rC5$9%5!<%P(B}
{A}
{$B8E$$%5!<%P$O%G%#%9%/>c32Ey$,$"$C$?$N$G!"<wL?$HH=CG$5$l!"%G!<%?$,B;<:$9$kA0$K?7$7$$%5!<%P$K0\9T$5$l$^$7$?!#(Bftp-master.debian.org$B$O(BDebian$B$N(B official package $B%j%]%8%H%j$G$9!#%Q%C%1!<%8$N%9%]%s%5!<(B(mentor)$B$rC5$9$N$O(Bmentors.debian.net$B!#(B }

\santaku
{pbuilder$B$K(Bclang support$B$,DI2C$5$l$^$7$?!#C/$,=q$$$?%Q%C%A$G$7$g$&$+!)(B}
{Sylvestre Ledru}
{Junichi Uekawa}
{Hideki Yamane}
{C}
{Debian$B$N(BClang$B%5%]!<%H$OCe!9$H?J$s$G$$$^$9!#(B}

\santaku
{DPN - 2013$BG/(B3$B7n(B4$BF|9f$K<h$j>e$2$i$l$?F|K\$N%$%Y%s%H$O(B}
{Open Source Conference 2013 Tokyo/Spring}
{Open Source Conference 2013 Hamamatu}
{Open Source Conference 2013 Tokushima}
{A}
{\url{http://henrich-on-debian.blogspot.jp/2013/02/open-source-conference-2013-tokyospring.html} $B>\:Y$O8e$[$I!#(B}


%\end{multicols}


% % (query-replace-regexp "<.*?>" "")
% % (query-replace-regexp "^[    ]\+" "")

%-------------------------------------------------------------------------------
\dancersection{Debian 10 busterでnftablesを使ってみる}{杉本典充}
%-------------------------------------------------------------------------------

\subsection{はじめに}

2019年7月6日にリリースした安定版Debian 10 busterではネットワークフィルタリング機能にnftablesがデフォルトで使われるように変わりました。

Debian 10 busterを使ってサーバを公開する場合にはnftablesを使うことがあると思いますので、調べてみました。


\subsection{nftableとは}

\subsubsection{nftableの歴史}

nftables\footnote{\url{https://wiki.nftables.org/wiki-nftables/index.php/Main_Page}}とは、linux-3.13 （2014/1/19リリース）から追加されたネットワークフィルタリング機能で iptables を置き換えることを目的として開発を進めている機能です。

iptables は linux-2.4 で新たに実装されたnetfilter\footnote{\url{https://www.netfilter.org/}}という機能を使ってネットワークフィルタリング処理を行う仕組みで、歴史が長いです\footnote{iptables-1.0.0.tar.bz2に含まれるmanファイル"iptables.8"には2000年3月20日という日付が入っています。}。

iptables は非常に長く多くの人たちに使われており、以下の課題があることがわかってきました\footnote{\url{https://wiki.nftables.org/wiki-nftables/index.php/Why_nftables\%3F}}。

\begin{itemize}
\item ソースコードに重複したコードが存在している
\item 性能の頭打ち
\item IPv4とIPv6のデュアルスタック時の管理の煩雑さ（nftablesのinetファミリーの追加につながる）
\item ルールセットの更新処理がアトミックでない
\item サードパーティアプリケーションが使えるAPIがない（nftablesではNetlink APIがある）
\item 構文の改善
\end{itemize}

上記を改善するため nftables の開発とlinux kernel の netfilter 機能の改善が行われ、2014年1月20日に nftables-0.099 がリリースされました\footnote{\url{https://www.netfilter.org/projects/nftables/downloads.html}}。現在の最新版は2019年12月2日にリリースした nftables-0.9.3 が最新版です。


\subsubsection{debianとnftables}

Debian においては、Debian wiki に nftables の情報があります\footnote{\url{https://wiki.debian.org/nftables}}。

Debian 9 stretch から nftables パッケージを提供しており、従来通り iptables をデフォルトとしつつ nftables も利用できる状況になっていました。このとき iptables が使う kernel module は以下になっています。

\begin{commandline}
$ cat /etc/debian_version
9.11  
$ lsmod | grep table
  ebtable_filter         16384  0
  ebtables               36864  1 ebtable_filter
  ip6table_filter        16384  0
  ip6_tables             28672  1 ip6table_filter
  iptable_filter         16384  1
  iptable_nat            16384  1
  nf_nat_ipv4            16384  1 iptable_nat
  iptable_mangle         16384  1
  ip_tables              24576  3 iptable_mangle,iptable_filter,iptable_nat
  x_tables               36864  12 ipt_REJECT,iptable_mangle,ip_tables,ebtables,iptable_filter,xt_tcpudp,
ipt_MASQUERADE,xt_CHECKSUM,ip6table_filter,xt_policy,xt_conntrack,ip6_tables
\end{commandline}

Debian 10 buster では iptables から nftables へデフォルトを切り替えており\footnote{\url{https://www.debian.org/releases/stable/amd64/release-notes/ch-whats-new.ja.html}}、nftables-0.9.0 を提供しています。nftables が使う kernel module は以下になっています。

\begin{commandline}
$ cat /etc/debian_version
10.2
$ lsmod | grep table
nf_tables_set          32768  6
nf_tables             143360  60 nft_ct,nf_tables_set
nfnetlink              16384  1 nf_tables
\end{commandline}

また、iptables パッケージの提供もしており昔ながらの使い方もできるようになっています。iptables パッケージの中には "/usr/sbin/iptables-nft" と "/usr/sbin/iptables-legacy" の二種類を提供しており、*-nft の方は nf\_tables の kernel module を使い、*-legacy の方は従来通り x\_tables の kernel module を使っています。そのため、Debian 10 busterでは、1. nf\_tables の nft インタフェース、2. nf\_tables の iptables 互換インタフェース、3. x\_tables の旧来の iptables インタフェース の3つが動作するようになっています。

Debian 11 bullseye で nftables と iptables をどうしていくかの意見がメーリングリストやブログに投稿されています。意見としては iptables パッケージを optional 扱いにして nftables パッケージのみインストールされるようにする案に賛同する人が見られます（他の意見は firewalld をデフォルトにしてはどうかという案）。

\begin{itemize}
\item default firewall utility changes for Debian 11 bullseye\footnote{\url{https://lists.debian.org/debian-devel/2019/07/msg00332.html}}
\item What to expect in Debian 11 Bullseye for nftables/iptables\footnote{\url{https://ral-arturo.org/2019/10/14/debian-netfilter.html}}
\end{itemize}


\subsection{Debian 10 busterでnftablesを使ってみる}

\subsubsection{インストール方法}

Debian 10 busterにおいて提供している nftables パッケージをインストールすると nft コマンドを利用できるようになり、nftables が systemd に登録されます。

\begin{commandline}
# apt-get install nftables
  
# which nft
/usr/sbin/nft
  
# systemctl status nftables
nftables.service - nftables
  Loaded: loaded (/lib/systemd/system/nftables.service; disabled; vendor preset: enabled)
  Active: inactive (dead)
  Docs: man:nft(8)
  http://wiki.nftables.org
\end{commandline}

nftables のパケットフィルタリングルールの設定ファイルは、/etc/nftables.conf に配置されます。初期状態では以下の設定が入っています。

\begin{commandline}
$ cat /etc/nftables.conf
#!/usr/sbin/nft -f

flush ruleset

table inet filter {
  chain input {
    type filter hook input priority 0;
  }
  chain forward {
    type filter hook forward priority 0;
  }
  chain output {
    type filter hook output priority 0;
  }
} 
\end{commandline}


\subsubsection{nftコマンドの使い方}

nftables によるネットワークフィルタリングの機能は nft コマンドで制御します。

"nft list ruleset" コマンドを実行すると現在のルールを表示します。

\begin{commandline}
# nft list ruleset
table inet filter {
    chain input {
        type filter hook input priority 0; policy accept;
    }

    chain forward {
        type filter hook forward priority 0; policy accept;
    }

    chain output {
        type filter hook output priority 0; policy accept;
    }
}
\end{commandline}

nftables へのルール定義の追加、変更、削除は以下の流れで行います。

\begin{itemize}
\item table の作成
\item chain の作成
\item rule の作成
\end{itemize}


\subsubsubsection{table の作成}

まずは最初に table を作成します。table を作成するときの引数にはアドレスファミリーを指定します。アドレスファミリーは以下表を値を指定できます\footnote{man nftより抜粋}。

\begin{table}[htb]
  \begin{center}
  \caption{nft で利用できるアドレスファミリー} 
  \begin{tabular}{|c|c|}
    \hline
    アドレスファミリー & 説明 \\ \hline
    ip & IPv4 address family. \\ \hline
    ip6 & IPv6 address family. \\ \hline
    inet & Internet (IPv4/IPv6) address family. \\ \hline
    arp & ARP address family, handling IPv4 ARP packets. \\ \hline
    bridge & Bridge address family, handling packets which traverse a bridge device. \\ \hline
    netdev & Netdev address family, handling packets from ingress. \\ \hline
  \end{tabular}
  \end{center}
\end{table}


コマンドの例です。

\begin{commandline}
# nft add table ip mytable
\end{commandline}


\subsubsubsection{chain の作成}

次に chain を作成します。chain を作成するときの引数には フックタイプ を指定します\footnote{linux kernelのNetfilter機構ではネットワークの入出力するパケットに対して任意な処理をひっかける機能がありこれを「フック」と呼んでいます。nftables や iptables で行う設定はこのフック処理を設定するフロントエンドツールと言えると思います。}。前述ではip (=IPv4）を指定していますので、以下表のIPv4/IPv6/Inetのフックタイプを指定できます\footnote{man nftより抜粋。}。

\begin{table}[htb]
  \begin{center}
  \caption{nft の IPv4/IPv6/Inet で利用できるフックタイプ} 
  \begin{tabular}{|c|c|}
    \hline
    フックタイプ & 説明 \\ \hline
    prerouting & ルーティング処理の前に実行するフック \\ \hline
    input & パケットの入力時に実行するフック \\ \hline
    forward & パケットの転送時に実行するフック \\ \hline
    output & パケットの出力時に実行するフック \\ \hline
    postrouting & ルーティング処理の後に呼び出されるフック \\ \hline
  \end{tabular}
  \end{center}
\end{table}

コマンドの例です。

\begin{commandline}
# nft add chain ip mytable mychain_in { type filter hook input priority 0 \; }
\end{commandline}


\subsubsubsection{rule の作成}

次に rule を作成をします。前述の例では IPv4 の入力パケットをフィルターするchainを作成しましたのでこれに準ずるruleを作成します。以下は、外部から入力パケット はssh のみを許可し、他は拒否する設定にしてみた例です。

\begin{commandline}
# nft add rule ip mytable mychain_in tcp dport ssh accept
# nft chain ip mytable mychain_in { policy drop \; }
\end{commandline}

上述の実行する nft コマンドを見ると「tcp」「dport」「ssh」「accept」と記述しています。「tcp」の部分は udp とも書け、「dport」の部分は sport とも書け、「ssh」の部分は プロトコル名またはポート番号でもよく、「accept」の部分は drop とも書けます。これらを記述を組み合わせて自分の実現したいルールを作ってください。

細かいruleを設定するコマンドの構文は「man nft」を参照してください。


\subsubsubsection{設定の保存}

ここまで設定すると以下のようなルールで動作しています\footnote{例示のためこのようなルールになっていますが、IPv6のssh通信が外からくると通信できてしまうため本番運用するにはおそらく設定として不完全ですので注意してください。}。

\begin{commandline}
# nft list ruleset
table inet filter {
    chain input {
        type filter hook input priority 0; policy accept;
    }

    chain forward {
        type filter hook forward priority 0; policy accept;
    }

    chain output {
        type filter hook output priority 0; policy accept;
    }
}
table ip mytable {
    chain mychain_in {
        type filter hook input priority 0; policy drop;
        tcp dport ssh accept
    }
}
\end{commandline}

nftables が参照する設定ファイルは /etc/nftables.conf であり、"nft list ruleset" コマンドの出力をそのまま保存すればよいようになっています。

\begin{commandline}
# cp -p /etc/nftables.conf /etc/nftables.conf.backup1
# nft list ruleset > /etc/nftables.conf  
\end{commandline}

nftables を再起動して同じ設定が読み込まれるか確認します。

\begin{commandline}
# systemctl restart nftables
# nft list ruleset
\end{commandline}

設定に問題がなければサーバの起動時にルールを自動反映するように設定します。

\begin{commandline}
# systemctl enable nftables
\end{commandline}


\subsection{利用シーンにおける設定例}

\subsubsection{インターネットに公開するwebサーバ向け}

インターネットに公開するwebサーバでは、基本的にhttp、https、sshのポートを公開すればよいと思います。設定は以下の条件で設定してみました。

\begin{itemize}
\item 入力パケットのデフォルト処理は破棄とし、許可したルールの入力パケットのみを受信
\item ループバックインタフェースの通信はすべて許可
\item 自分のホストから通信を開始した場合の相手のサーバからの戻りパケットは受信許可
\item TCPのコネクションの状態としておかしいパケットは破棄
\item pingのエコー要求を許可
\item sshをLAN内から通信を許可
\item sshをインターネットからポート番号をtcpの10022ポートに変更した状態で通信を許可
\item http/httpsはインターネットからもLAN内からも通信を許可
\item IPv6ではssh、http/httpsのサービスはせず、ping6の応答もしない
\end{itemize}

nftables の設定は以下になります。

\begin{commandline}
# nft list ruleset
table ip filter {
    chain input {
        type filter hook input priority 0; policy drop;
        ct state { established, related } accept
        ct state { invalid } drop
        iifname "lo" accept
        icmp type { echo-reply, echo-request } accept
        tcp dport ssh ip saddr { 10.0.0.0/8, 172.16.0.0/12, 192.168.0.0/16 } accept
        tcp dport 10022 ip saddr { 0.0.0.0-255.255.255.255 } accept
        tcp dport { http, https } ip saddr { 0.0.0.0-255.255.255.255 } accept
    }

    chain forward {
        type filter hook forward priority 0; policy accept;
    }

    chain output {
        type filter hook output priority 0; policy accept;
    }
}
table ip6 filter {
    chain input {
        type filter hook input priority 0; policy drop;
        iifname "lo" accept
    }

    chain forward {
        type filter hook forward priority 0; policy accept;
    }

    chain output {
        type filter hook output priority 0; policy accept;
    }
}
\end{commandline}


\subsubsection{インターネットに公開するopenvpnサーバ向け}

openvpn\footnote{\url{https://openvpn.net/community/}} とは SSL-VPN の機能をもった VPN サーバ及び VPN クライアントアプリケーションです。openvpn を使って VPN 接続したクライアントは openvpn サーバを中継して別のネットワークのサーバと通信することができます。

ここでは openvpn に接続したサーバ及びクライアントに割り当てられるネットワークアドレス帯"192.168.200.0/24"からパケットが入力された場合に、他のネットワークへパケットをルーティングする設定を行います。

\begin{commandline}
# nft add table ip nat
# nft add chain ip nat postrouting { type nat hook postrouting priority 100 \; }
# nft add rule ip nat postrouting oifname "eth0" ip saddr 192.168.200.0/24 masquerade 
\end{commandline}

openvpnを起動すると追加される"tun"インタフェースについても受信許可し、openvpnが利用するUDPとTCPの119番ポートの通信を受信許可に設定します。

nftables の設定は以下になります。

\begin{commandline}
# nft list ruleset
table ip filter {
    chain input {
        type filter hook input priority 0; policy drop;
        ct state { established, related } accept
        ct state { invalid } drop
        iifname "lo" accept
        iifname "tun*" accept
        icmp type { echo-reply, echo-request } accept
        tcp dport ssh ip saddr { 10.0.0.0/8, 172.16.0.0/12, 192.168.0.0/16 } accept
        tcp dport 10022 ip saddr { 0.0.0.0-255.255.255.255 } accept
        udp dport openvpn ip saddr { 0.0.0.0-255.255.255.255 } accept
        tcp dport openvpn ip saddr { 0.0.0.0-255.255.255.255 } accept
    }

    chain forward {
        type filter hook forward priority 0; policy accept;
    }

    chain output {
        type filter hook output priority 0; policy accept;
    }
}
table ip6 filter {
    chain input {
        type filter hook input priority 0; policy drop;
        iifname "lo" accept
    }

    chain forward {
        type filter hook forward priority 0; policy accept;
    }

    chain output {
        type filter hook output priority 0; policy accept;
    }
}  
table ip nat {
    chain postrouting {
        type nat hook postrouting priority 100; policy accept;
        oifname "eth0" ip saddr 192.168.200.0/24 masquerade
    }
}
\end{commandline}


\subsection{おわりに}

Debian 10 でデフォルトになった nftables を調べてみました。

今後 nftables が主流になっていくと思われるため Debian 10 buster を使い始める方は nftables も合わせて使い始めてみてはいかがでしょうか。 


\subsection{参考文献}

\begin{itemize}
\item 「nftables - Debian Wiki」 \url{https://wiki.debian.org/nftables}
\item 「netfilte.org」 \url{https://www.netfilter.org/}
\item 「netfilter wiki」 \url{https://wiki.nftables.org/wiki-nftables/index.php/Main_Page}
\item 「Linuxにおける新たなパケットフィルタリングツール「nftables」入門 - さくらのナレッジ」 \url{https://knowledge.sakura.ad.jp/22636/}
\item 「nftables で基本的なフィルタリングを設定してみた - SIOS TECH.LAB」 \url{https://tech-lab.sios.jp/archives/16930}
\end{itemize}


% 冊子にするために、4の倍数にする必要がある。
% そのための調整
%%\dancersection{メモ}{}
%%\mbox{}\newpage

\vspace*{15cm}
\hrule
\vspace{2mm}
\includegraphics[width=2cm]{image200502/openlogo-nd.eps}
\noindent \Large \bf Debian 勉強会資料\\
\noindent \normalfont \debmtgyear{}年\debmtgmonth{}月\debmtgdate{}日 \hspace{5mm}  初版第1刷発行\\
\noindent \normalfont 東京エリア Debian 勉強会 （編集・印刷・発行）\\
\hrule

\end{document}
